% ============================================================
%  SECCIÓN EXTENDIDA — Para insertar DESPUÉS de \section{Conclusiones}
%  y ANTES de \printbibliography en paper_TEA_AACC_bibliometrico.tex
%
%  Instrucciones de integración en Overleaf:
%  1. Abrir paper_TEA_AACC_bibliometrico.tex
%  2. Localizar la línea:  \newpage\printbibliography[title={Referencias}]
%  3. Pegar todo el contenido de este archivo ANTES de esa línea
%  4. Añadir las entradas de referencias.bib adicionales que aparecen
%     al final de este archivo
% ============================================================

% ═══════════════════════════════════════════════════════════
\section{Revisión crítica de los trabajos existentes}
\label{sec:critica}
% ═══════════════════════════════════════════════════════════

El reducido corpus de estudios empíricos específicos sobre TEA+AACC en contextos
universitarios y de investigación —identificado en la presente revisión en apenas
24 trabajos con mención explícita de educación superior o adultos profesionales—
permite un análisis crítico detallado de sus aportes, limitaciones y tensiones
metodológicas. A continuación, se examinan los estudios más relevantes agrupados
por línea temática.

\subsection{Estudios sobre perfiles cognitivos y trayectorias académicas}

El trabajo fundacional de \textcite{foleynicpon2011_paradox} estableció el marco
conceptual de la ``paradoja de la superdotación y el autismo'': la misma
arquitectura cognitiva que genera capacidades excepcionales —el procesamiento
local hiperfocalizado, la atención sostenida a los detalles, la memoria semántica
profunda— puede, en determinadas condiciones, traducirse en las dificultades de
flexibilidad cognitiva, coherencia central y regulación ejecutiva que caracterizan
al TEA. Esta tensión no es una contradicción sino una \textit{bifurcación
expresiva} de un mismo sustrato neurológico.

\textcite{cain2019} ofrecen, hasta la fecha, el estudio más riguroso sobre
trayectorias académicas de niños con CI$\geq$130 y diagnóstico confirmado de TEA.
Con una muestra de 71 participantes seguidos durante tres años, los autores
documentan que el grupo 2eASD presenta rendimiento académico significativamente
inferior al esperado por su capacidad intelectual en todas las asignaturas que
implican procesamiento social o comunicación escrita (composición, ciencias
sociales), mientras que su rendimiento en matemáticas y ciencias naturales
no difiere significativamente del grupo dotado sin TEA. Esta \textit{disociación
de dominio} tiene implicaciones directas para el diseño curricular en educación
universitaria y para la evaluación del rendimiento investigador.

\textbf{Crítica metodológica:} La muestra de \textcite{cain2019} proviene
íntegramente de clínicas especializadas en trastornos del neurodesarrollo en
California (EE.UU.), introduciendo un sesgo de selección clínica severo: la
mayoría de los adultos 2e en entornos académicos nunca han accedido a evaluación
clínica especializada y, por tanto, quedan sistemáticamente excluidos de este
tipo de estudios. Adicionalmente, el seguimiento de tres años solo cubre la
etapa de educación primaria, dejando abierta la pregunta central sobre qué
ocurre en la transición a educación superior y en la carrera profesional.

\textcite{assouline2021} aportan evidencia sobre el papel mediador de la
\textit{velocidad de procesamiento} (VP) en el rendimiento académico de
estudiantes 2eASD. Sus análisis de regresión múltiple sobre una muestra de 152
estudiantes de secundaria muestran que la VP explica el 34\,\% de la varianza
en el rendimiento académico, por encima de la capacidad verbal y el razonamiento
fluido. Este hallazgo es especialmente relevante para el contexto universitario
e investigador: las evaluaciones académicas con límite de tiempo, los exámenes
orales sin preparación previa, las defensas de tesis y la redacción bajo presión
son todos escenarios que penalizan sistemáticamente a individuos con VP baja
independientemente de su capacidad real.

\subsection{Estudios en educación superior: un corpus mínimo pero esclarecedor}

Solo se identificaron 24 trabajos con foco en educación universitaria en el
corpus completo, de los cuales apenas 7 abordan específicamente a estudiantes
o profesionales con la combinación TEA+AACC (doble excepcionalidad), y ninguno
estudia a académicos como \textit{sujetos de investigación} en lugar de como
\textit{facilitadores} del proceso educativo de otros.

\subsubsection{La triada Reis-Gelbar-Madaus y el modelo UConn}

El grupo de investigación de la Universidad de Connecticut —liderado por Reis,
Gelbar y Madaus— ha producido la serie más sistemática de estudios sobre
estudiantes 2eASD en educación superior
\parencite{reis2021, madaus2022factors, madaus2022transition, austermann2023}.
Sus hallazgos principales merecen una lectura crítica detallada:

\textcite{reis2021} analizan mediante estudio de casos múltiples a 40 estudiantes
de colegios universitarios altamente competitivos diagnosticados con TEA y con
antecedentes de alto rendimiento académico. El estudio identifica tres factores
como predictores del éxito académico: (1) acceso temprano a intervención
conductual, (2) entornos familiares con altas expectativas y apoyo explícito a
la neurodiversidad, y (3) presencia de al menos un mentor académico que reconozca
y valore el perfil neurodivergente. La \textit{autocomprensión del perfil
neurodivergente} —entender sus propias fortalezas y dificultades— correlaciona
positivamente con el rendimiento universitario independientemente del CI.

\textbf{Crítica:} La muestra de Reis et al.\ se limita a universidades de élite
norteamericanas con servicios de apoyo a la discapacidad bien desarrollados,
lo que crea un sesgo de supervivencia doble: solo los estudiantes 2eASD con
mayores recursos (económicos, familiares, terapéuticos) llegan a estas
instituciones y persisten en ellas. Los hallazgos son difícilmente transferibles
a contextos con menor dotación de recursos, como las universidades públicas
latinoamericanas.

\textcite{madaus2022factors} aplican un diseño de grupos de enfoque con 23 estudiantes
2eASD de varias universidades norteamericanas para identificar factores
facilitadores e inhibidores del aprendizaje. Los obstáculos más frecuentemente
citados son: la ambigüedad en las instrucciones de los profesores (87\,\% de los
participantes), los trabajos grupales no estructurados (78\,\%), la evaluación
de presentaciones orales (73\,\%), y los campus ruidosos y sensorialmente
sobreestimulantes (69\,\%). Los facilitadores más valorados son la claridad
estructural, el acceso a materiales por adelantado y la posibilidad de entregar
trabajos en formato alternativo.

\subsubsection{Estudios cualitativos en el Reino Unido}

\textcite{macleod2017} realizan un estudio cualitativo longitudinal (18 meses)
con 12 estudiantes autistas en universidades escocesas. El hallazgo central
es lo que los autores denominan ``el coste del éxito'': los estudiantes más
capaces —precisamente aquellos que logran altas calificaciones— son los que
acumulan mayor agotamiento psicológico, porque el mantenimiento del rendimiento
exige un esfuerzo cognitivo y social desproporcionadamente mayor que para sus
pares neurotípicos. Este fenómeno, que los autores equiparan clínicamente al
burnout, es descrito por los participantes como ``actuar constantemente en una
obra de teatro para la que no me dieron el guión''.

\textbf{Crítica metodológica:} La muestra de \textcite{macleod2017} no
distingue entre estudiantes con y sin altas capacidades, lo que impide
determinar si el ``coste del éxito'' es más pronunciado en perfiles 2e. Esta
es una limitación recurrente en la literatura: la mayoría de los estudios
sobre autismo en educación superior no controlan por capacidad intelectual,
lo que equivale a tratar el espectro como homogéneo cuando la heterogeneidad
es una de sus características definitivas.

\textcite{gillespie-lynch2017} amplían la perspectiva incluyendo las voces de
los propios estudiantes neurodivergentes en el diseño de apoyos universitarios.
Su metodología de investigación participativa —en la que los estudiantes son
co-investigadores— es metodológicamente innovadora y éticamente relevante:
produce recomendaciones con mayor validez ecológica que los estudios diseñados
exclusivamente por investigadores neurotípicos. Los autores documentan que las
modificaciones curriculares más efectivas son frecuentemente de \textit{bajo
coste} (instrucciones escritas en lugar de orales, rúbricas detalladas,
feedbacks estructurados) pero requieren un cambio en la \textit{cultura
institucional}, no solo en los procedimientos formales.

\subsubsection{El estudio de Perkowski y el caso del doctorado}

El trabajo de \textcite{perkowski2017} representa el único estudio identificado
en el corpus que aborda explícitamente la relación supervisor--doctorando cuando
el estudiante tiene diagnóstico de TEA. Mediante un estudio de caso etnográfico
en una universidad polaca, los autores documentan cómo la incomprensión del
perfil autístico por parte del director de tesis genera dinámicas de conflicto
que tienen poco que ver con la capacidad intelectual del doctorando y mucho con
la incompatibilidad entre los estilos comunicativos y las expectativas relacionales
implícitas en la relación académica tradicional. El doctorando del caso, que
posteriormente publicó investigación de alta calidad, estuvo a punto de abandonar
el programa doctoral exclusivamente por estas fricciones relacionales.

Este estudio, aunque limitado por su $n=1$ y su contexto geográfico específico,
ilustra un mecanismo de pérdida de talento investigador que probablemente se
reproduce con alta frecuencia en sistemas académicos de todo el mundo, incluyendo
el colombiano, y que permanece completamente invisible en las estadísticas de
abandono doctoral.

\subsection{Instrumentos de evaluación: avances y limitaciones}

\textcite{cederberg2018} evaluaron la validez de dos instrumentos de cribado del
TEA —el ASSQ y el SRS— en muestras de jóvenes con alta capacidad intelectual,
encontrando que ambos instrumentos presentan tasas de falsos negativos
significativamente más altas en personas con CI$\geq$130 que en la población
general. El enmascaramiento cognitivo opera directamente sobre los ítems de estos
cuestionarios: las personas con AACC aprenden a dar las respuestas ``esperadas''
en contextos de evaluación.

\textcite{alhroub2021} propone el uso de la \textit{evaluación dinámica}
—basada en la Zona de Desarrollo Próximo vygotskiana— como complemento a las
pruebas psicométricas estáticas para identificar el perfil 2e. La evaluación
dinámica mide la capacidad de aprendizaje mediado, no el rendimiento actual,
lo que puede revelar potencial en individuos cuyo desempeño actual está
suprimido por las dificultades asociadas al TEA. Esta propuesta es
metodológicamente prometedora pero carece aún de validación en muestras de
adultos universitarios.

\textcite{burgerveltmeijer2023} presentan la heurística S\&W (Strengths \&
Weaknesses), un protocolo de evaluación multidimensional para estudiantes 2e
que combina: (1) evaluación psicométrica del perfil cognitivo; (2) entrevista
semiestructurada de historia de vida; (3) observación en contexto académico; y
(4) cuestionario de autoinforme de estrategias de afrontamiento. Esta heurística,
desarrollada y validada en los Países Bajos, representa el avance metodológico
más completo hasta la fecha para la evaluación del perfil 2eASD, pero no ha
sido adaptada al contexto latinoamericano ni validada en muestras de adultos
universitarios o investigadores.

\subsection{El bienestar psicológico del adulto 2e: un campo prácticamente virgen}

\textcite{wallpolin2017} realizan, en lo que nuestro corpus revela como el único
estudio específicamente centrado en la calidad de vida de \textit{adultos}
2eASD, entrevistas en profundidad con 15 adultos con diagnóstico de TEA e
historia documentada de altas capacidades. El título del estudio —``Relative
to Einstein'' (En relación a Einstein)— captura la ambivalencia central del
perfil: los participantes oscilan entre una narrativa de orgullo por sus
capacidades excepcionales y una narrativa de agotamiento por el esfuerzo
permanente de adaptación a un mundo diseñado para el neurotipo mayoritario.
La mayoría describían sus vidas profesionales como exitosas \textit{en
apariencia} pero profundamente insatisfactorias en términos de costo personal.

\textcite{poirier2025} amplían este análisis con un diseño cuantitativo-cualitativo
en 68 adultos con AACC (con y sin diagnóstico de TEA), encontrando que la
presencia de un diagnóstico de TEA modera significativamente la relación entre
alta capacidad intelectual y bienestar subjetivo: mientras que los adultos
dotados sin TEA reportan niveles de bienestar superiores a la media, los adultos
2e reportan niveles inferiores, con especial impacto en las dimensiones de
satisfacción relacional, sentido de pertenencia y coherencia narrativa de la
identidad.

\textbf{Vacío crítico:} No existe en el corpus ningún estudio que examine
específicamente el bienestar psicológico de académicos o investigadores
profesionales con TEA y AACC en el contexto de sus carreras científicas,
con la sola excepción parcial del trabajo de \textcite{macleod2017} centrado
en estudiantes universitarios. La dimensión específicamente profesional
—la relación entre el perfil 2e y la productividad investigadora, el sistema
de evaluación meritocrática académica, la precariedad laboral del sector,
o la cultura de hiperproductividad de la academia contemporánea— permanece
completamente inexplorada.

\subsection{Tensiones conceptuales y debates abiertos en el campo}

\subsubsection{¿Neurodiversidad como fortaleza o como vulnerabilidad?}

Un debate fundamental atraviesa la literatura y raramente se tematiza de forma
explícita: si el perfil 2eASD debe conceptualizarse como una \textit{condición}
que requiere apoyo terapéutico y compensatorio, o como una \textit{variante
neurodivergente} que requiere entornos adaptativos pero no corrección. Esta
tensión no es meramente terminológica: condiciona los instrumentos utilizados,
las hipótesis formuladas, los resultados considerados relevantes y, en última
instancia, las políticas recomendadas.

\textcite{linares2026} proponen un modelo integrativo entre el paradigma
clínico-psicológico (centrado en déficits, apoyos y calidad de vida) y el
paradigma de la neurodiversidad (centrado en diferencia, identidad y entorno
adaptativo). Su argumento es que ambas perspectivas son necesarias y que su
integración —aún incipiente en la literatura— permitiría diseñar intervenciones
más respetuosas de la agencia de las personas 2e. Este marco integrativo
tiene especial relevancia para el diseño de programas de apoyo en universidades,
donde coexisten lógicas institucionales de gestión de la discapacidad
(certificaciones, acomodaciones, servicios de bienestar) con comunidades
emergentes de activismo neurodivergente.

\subsubsection{El diagnóstico tardío y la reconfiguración identitaria}

\textcite{pessoa2026}, en un trabajo reciente desde Brasil, documentan el proceso
de reconfiguración identitaria que experimentan adolescentes y adultos con AACC
y TEA que reciben el diagnóstico de forma tardía. Los autores identifican cuatro
fases: (1) incredulidad y resistencia (``yo no soy así''); (2) reconocimiento
retrospectivo (``todo encaja ahora''); (3) duelo por el tiempo perdido sin apoyo;
y (4) integración narrativa y empoderamiento. Este proceso, documentado en el
contexto brasileño, probablemente presenta variaciones culturales relevantes
en el contexto colombiano, donde el estigma asociado al diagnóstico psiquiátrico
y la exaltación del rendimiento académico como marcador de éxito social
interactuarían con este proceso de formas específicas que merecen investigación.

\subsubsection{Crítica a la noción de ``éxito académico'' como variable de resultado}

Varios autores \parencite{macleod2017, gillespie-lynch2017, wallpolin2017}
cuestionan la utilización del rendimiento académico y la productividad
como las principales variables de resultado en los estudios sobre 2eASD.
Esta crítica es de fondo: medir el ``éxito'' de las personas 2e únicamente
por su capacidad de adaptarse a sistemas evaluativos diseñados para el
neurotipo mayoritario equivale a preguntar si un daltónico ``tiene éxito''
en distinguir semáforos sin adaptar el entorno. Lo que los estudios
identifican como éxito puede ser, en muchos casos, la manifestación de un
enmascaramiento costoso y sostenido cuya factura se paga en forma de
burnout, crisis de salud mental y abandono prematuro de carreras
prometedoras.

% ═══════════════════════════════════════════════════════════
\section{Metodologías propuestas para el estudio del fenómeno en contexto
         universitario: hacia un programa de investigación}
\label{sec:metodologias}
% ═══════════════════════════════════════════════════════════

La revisión anterior permite identificar con precisión no solo los vacíos
de conocimiento sino también las exigencias metodológicas que un programa
de investigación sobre TEA+AACC en académicos y profesionales debería cumplir.
Se propone un marco de siete aproximaciones complementarias, ordenadas de menor
a mayor complejidad, que pueden implementarse progresivamente y que son
especialmente pertinentes para el contexto de universidades latinoamericanas con
recursos limitados.

\subsection{Fase 1 — Exploración y detección: estudios de prevalencia y cribado}

\subsubsection{Diseño}

Estudio transversal descriptivo de prevalencia del perfil 2e en comunidades
universitarias (docentes, investigadores, estudiantes de posgrado), utilizando
instrumentos de cribado validados en español.

\subsubsection{Instrumentos propuestos}

\begin{enumerate}[leftmargin=*, label=\textbf{\arabic*.}, itemsep=3pt]
    \item \textbf{Autism Quotient (AQ-28 versión española, \textcite{baron-cohen2001})}.
    Cuestionario de autoinforme de 28 ítems con validación en población hispanohablante.
    Sensibilidad del 72\,\% y especificidad del 82\,\% para detección de rasgos
    autísticos en población general. Limitación: alta tasa de falsos negativos en
    personas con AACC por capacidad de respuesta estratégica.

    \item \textbf{Ritvo Autism Asperger Diagnostic Scale-Revised (RAADS-R)}.
    80 ítems, diseñado específicamente para adultos de alta funcionalidad con
    sospecha de TEA. Mayor sensibilidad que el AQ en perfiles compensados.
    Disponible en versión española.

    \item \textbf{Gifted Rating Scales (GRS-S / GRS-T) adaptadas}.
    O, alternativamente, los criterios del \textit{Columbus Group} (1991) para
    identificación de altas capacidades en adultos. Para el contexto colombiano,
    se recomienda complementar con el modelo de Renzulli adaptado por
    \textcite{gonzalez2010} para población latinoamericana.

    \item \textbf{Camouflaging Autistic Traits Questionnaire (CAT-Q)}.
    25 ítems que miden específicamente las estrategias de enmascaramiento
    (\textit{masking}), con tres subescalas: asimilación, compensación y enmascaramiento.
    Especialmente relevante para población universitaria donde el enmascaramiento
    es pronunciado. Disponible en español.
\end{enumerate}

\subsubsection{Procedimiento}

Aplicación online anónima precedida por información detallada sobre el propósito
del estudio y los conceptos de doble excepcionalidad. El anonimato es condición
necesaria para reducir la inhibición de respuesta por estigma. Los participantes
con puntaje sugestivo deben ser informados de los recursos de evaluación
diagnóstica disponibles, sin que el estudio constituya en sí mismo un diagnóstico.

\subsubsection{Consideraciones éticas}

La investigación con poblaciones universitarias implica riesgos de
re-estigmatización y debe cumplir estrictamente los principios de:
(a) consentimiento informado explícito; (b) anonimato garantizado técnicamente
(no solo prometido); (c) ausencia de consecuencias laborales o académicas
derivadas de la participación; y (d) disponibilidad de apoyo psicológico
para participantes que experimenten malestar al responder los cuestionarios.

\subsection{Fase 2 — Comprensión: estudios cualitativos de experiencia vivida}

\subsubsection{Diseño}

Investigación fenomenológica interpretativa (IPA, \textit{Interpretative
Phenomenological Analysis}, \textcite{smith2009}) o narrativa con docentes,
investigadores y estudiantes de posgrado con diagnóstico de TEA (con o sin
AACC documentadas).

\subsubsection{Guía de preguntas}

Las entrevistas en profundidad deben explorar: (1) trayectoria diagnóstica
(cuándo, cómo, qué cambió con el diagnóstico); (2) experiencias en entornos
académicos (docencia, investigación, publicación, evaluación de pares,
networking); (3) estrategias de afrontamiento y enmascaramiento utilizadas;
(4) relación con pares, supervisores y subordinados; (5) percepción de la
propia capacidad intelectual en relación al TEA; (6) necesidades de apoyo
no satisfechas por la institución.

\subsubsection{Ventaja metodológica para contextos de recursos limitados}

La IPA y los estudios narrativos requieren muestras pequeñas ($n=6$--$15$)
y producen conocimiento profundo y contextualy transferible. Son metodologías
adecuadas para una primera fase de investigación en instituciones con recursos
limitados para financiar estudios de gran escala.

\subsection{Fase 3 — Medición: estudios de caso de perfil neuropsicológico}

\subsubsection{Diseño}

Estudio de caso múltiple con evaluación neuropsicológica comprehensiva aplicando
la heurística S\&W \parencite{burgerveltmeijer2023} adaptada al contexto adulto:

\begin{enumerate}[leftmargin=*, itemsep=2pt]
    \item \textbf{Evaluación del perfil cognitivo}: WAIS-IV (o RIAS-2),
    con análisis de índices específicos (Índice de Velocidad de Procesamiento,
    Índice de Memoria de Trabajo, Índice de Comprensión Verbal, Índice de
    Razonamiento Perceptivo). El patrón típico 2eASD presenta ICV muy elevado
    (P$\geq$95) con IVP marcadamente inferior.
    \item \textbf{Evaluación de funciones ejecutivas}: BRIEF-A (Behavior Rating
    Inventory of Executive Function -- Adult version).
    \item \textbf{Evaluación de rasgos autísticos}: ADOS-2 módulo 4 para adultos
    de alta funcionalidad, complementado con ADI-R.
    \item \textbf{Historia de vida académica}: entrevista semiestructurada
    específica para el perfil investigador-académico.
    \item \textbf{Análisis de producción científica}: CV, publicaciones, proyectos
    de investigación, participación en redes académicas.
\end{enumerate}

\subsubsection{Aporte específico al contexto colombiano}

Este diseño permite documentar con precisión el perfil de académicos colombianos
2e, aportando datos normativos locales que actualmente no existen. El estudio
de caso no pretende generalización estadística sino profundidad analítica:
la comprensión de uno o dos casos en detalle puede generar hipótesis y teorías
transferibles que guíen investigaciones cuantitativas posteriores
\parencite{yin2018}.

\subsection{Fase 4 — Acción: investigación-acción participativa (IAP)}

\subsubsection{Diseño}

Modelo de IAP \parencite{kemmis2014} en el que los propios académicos 2e participan
en el diseño, implementación y evaluación de programas de apoyo institucional.
Los ciclos de acción-reflexión permiten adaptar las intervenciones a las
necesidades reales de los participantes, superando la limitación de los
programas importados de contextos anglosajones sin adaptación cultural.

\subsubsection{Componentes de la intervención}

\begin{itemize}[leftmargin=*, itemsep=2pt]
    \item \textbf{Grupos de pares neurodivergentes} (peer support groups):
    espacios de intercambio de experiencias y estrategias entre docentes e
    investigadores 2e. La evidencia sobre el valor del apoyo entre pares
    neurodivergentes es consistente \parencite{gillespie-lynch2017}.
    \item \textbf{Formación a coordinadores académicos y pares}: talleres
    sobre perfil 2e, reconocimiento de señales de dificultad enmascarada,
    ajustes razonables no estigmatizantes.
    \item \textbf{Protocolo de ajustes razonables para investigadores}: adaptación
    de cargos de trabajo, plazos de entrega, formatos de evaluación y condiciones
    del ambiente físico (luz, ruido, temperatura) para el personal académico con
    TEA. Este componente requiere coordinación con las áreas de Talento Humano
    y Bienestar Universitario.
\end{itemize}

\subsection{Fase 5 — Validación: estudios comparativos multicéntricos}

Una vez establecidas bases de datos en universidades individuales, el siguiente
paso metodológico es la comparación entre instituciones con diferentes
características (tamaño, tipo, ubicación geográfica, oferta académica) para
identificar factores institucionales que modulan la experiencia y el rendimiento
de los académicos 2e. En Colombia, una comparación entre universidades públicas
de distintos tamaños y regiones —Unisucre en Sincelejo, Universidad de los Andes
en Bogotá, Universidad de Antioquia en Medellín, Universidad del Pacífico en
Buenaventura— permitiría controlar por factores socioeconómicos, culturales y
de dotación de recursos que el foco exclusivo en élites académicas no puede capturar.

% ═══════════════════════════════════════════════════════════
\section{Aplicación en el contexto de una universidad pequeña colombiana:
         el caso de Unisucre}
\label{sec:unisucre}
% ═══════════════════════════════════════════════════════════

\subsection{Caracterización del contexto institucional}

La Universidad de Sucre (Unisucre), institución pública de educación superior con
sede en Sincelejo (Sucre, Colombia), representa un caso de estudio de particular
pertinencia para la agenda de investigación propuesta. Con aproximadamente 8.000
estudiantes, 12 programas académicos de pregrado y una planta docente de alrededor
de 200 profesores de tiempo completo, Unisucre opera en un contexto geográfico
y socioeconómico con características que intensifican los desafíos propios del
perfil 2e: alta prevalencia de vulnerabilidad socioeconómica en la comunidad
estudiantil, recursos limitados para servicios de bienestar y salud mental,
ausencia histórica de programas de atención a la neurodiversidad, y una cultura
institucional en proceso de transición hacia modelos de investigación y calidad
académica exigentes.

Al mismo tiempo, Unisucre ofrece condiciones favorables para una investigación
pionera: su tamaño manejable facilita el acceso a la comunidad docente e
investigadora como población de estudio, su proceso de acreditación institucional
genera incentivos para documentar y mejorar prácticas de bienestar y equidad, y
la inexistencia de investigación previa convierte a cualquier trabajo en el área
en una contribución de primera línea al conocimiento.

\subsection{Estimación del tamaño de la población potencial}

Aunque no existen datos específicos, es posible hacer una estimación conservadora
basada en la evidencia disponible. Si la prevalencia de TEA en adultos de alta
funcionalidad se estima en torno al 1--2\,\% de la población general
\parencite{maenner2023}, y si la proporción de académicos con perfil 2e (TEA+AACC)
es consistente con la mayor representación de rasgos autísticos en entornos
STEM documentada por \textcite{baron-cohen2001}, cabría esperar que entre el
2\,\% y el 5\,\% de la planta docente de una universidad presentara un perfil
2e no diagnosticado o con diagnóstico tardío. En Unisucre, esto representaría
entre 4 y 10 docentes en activo —una población perfectamente manejable para
un estudio de caso cualitativo de fase 1.

Esta estimación es deliberadamente conservadora y puede estar subestimando
la prevalencia real por dos razones: (1) el subdiagnóstico es especialmente
pronunciado en Colombia por las barreras de acceso a evaluación especializada;
y (2) el perfil investigador-académico puede atraer desproporcionadamente a
personas con rasgos 2e, como sugiere la evidencia internacional.

\subsection{Propuesta de protocolo adaptado a Unisucre}

\subsubsection{Fase 0 — Sensibilización institucional y alianzas (meses 1--3)}

Antes de iniciar cualquier recolección de datos, es indispensable construir
alianzas institucionales que garanticen la sostenibilidad del proceso y la
protección de los participantes:

\begin{itemize}[leftmargin=*, itemsep=2pt]
    \item \textbf{Comité de Ética}: obtención de aval del Comité Institucional
    de Ética de Investigación de Unisucre, siguiendo la Resolución 8430 de 1993
    del Ministerio de Salud colombiano (clasificación de riesgo mínimo).
    \item \textbf{Bienestar Universitario}: articulación con la Unidad de Bienestar
    para garantizar acompañamiento psicológico disponible para participantes
    que lo requieran.
    \item \textbf{Vicerrectoría Académica}: compromisos de confidencialidad y
    neutralidad institucional — la participación (o no participación) no puede
    tener consecuencias sobre la vinculación o la evaluación docente.
    \item \textbf{Grupos estudiantiles}: especialmente estudiantes de psicología,
    trabajo social y ciencias de la salud como potenciales colaboradores del proyecto.
\end{itemize}

\subsubsection{Fase 1 — Cribado (meses 4--8)}

\begin{itemize}[leftmargin=*, itemsep=2pt]
    \item Convocatoria voluntaria y anónima a toda la planta docente mediante
    correo institucional, con explicación del propósito del estudio y los conceptos
    de doble excepcionalidad y neurodiversidad.
    \item Aplicación de batería de cribado online: AQ-28 + CAT-Q + indicadores
    de alta capacidad (AQ-CI adaptado o criterios Columbus Group).
    \item Umbral de detección conservador (percentil 70 en AQ-28 más cualquier
    puntuación alta en CAT-Q) para maximizar sensibilidad dado el alto
    enmascaramiento esperado.
    \item Los participantes reciben retroalimentación individual automatizada sobre
    sus puntuaciones, con recursos informativos sobre TEA y AACC y oferta de
    participación voluntaria en la siguiente fase.
\end{itemize}

\subsubsection{Fase 2 — Entrevistas en profundidad (meses 9--14)}

Con los docentes que superen el umbral de cribado y consientan participar
($n$ esperado: 6--12), se realizan entrevistas semiestructuradas individuales
de 90--120 minutos con los siguientes bloques temáticos:

\begin{enumerate}[leftmargin=*, label=\textbf{B\arabic*.}, itemsep=3pt]
    \item \textbf{Trayectoria biográfica y académica}: reconstrucción narrativa
    del recorrido educativo y profesional, con atención especial a momentos de
    ``no encaje'' y estrategias de adaptación.
    \item \textbf{Experiencia docente e investigadora}: cómo se vive la docencia,
    la investigación, la publicación, la evaluación por pares y las interacciones
    con colegas y estudiantes.
    \item \textbf{Diagnóstico (si existe) y sus efectos}: cuándo, cómo, qué
    cambió. Para quienes no tienen diagnóstico: experiencia de autoidentificación.
    \item \textbf{Necesidades de apoyo}: qué apoyos institucionales serían
    significativos y qué barreras existen actualmente.
    \item \textbf{Contexto cultural y regional}: cómo influyen la cultura costeña,
    las expectativas familiares y comunitarias, el estigma local hacia los
    diagnósticos psiquiátricos.
\end{enumerate}

\subsubsection{Fase 3 — Validación e impacto (meses 15--24)}

\begin{itemize}[leftmargin=*, itemsep=2pt]
    \item Análisis temático de las entrevistas (método de Braun y Clarke, 2006).
    \item Devolución de resultados preliminares a los participantes para validación
    (\textit{member checking}).
    \item Co-diseño con los participantes de un \textbf{protocolo de ajustes
    razonables} adaptado al contexto de Unisucre, que incluya recomendaciones
    específicas para: (a) carga docente y asignación de grupos; (b) participación
    en comités y actividades de extensión; (c) evaluación docente; (d) condiciones
    del ambiente físico de trabajo.
    \item Diseminación: publicación en revistas colombianas indexadas (Revista
    Colombiana de Educación, Universitas Psychologica) + presentación en Coloquios
    Regionales de Educación Inclusiva.
\end{itemize}

\subsection{Retos específicos del contexto de Unisucre}

\subsubsection{Barrera diagnóstica y terminológica}

En Sincelejo y la Costa Caribe colombiana, el acceso a evaluación neuropsicológica
especializada es extremadamente limitado. Los tres psiquiatras o psicólogos
clínicos con experiencia en TEA adulto más cercanos a Unisucre se encuentran en
Barranquilla o Cartagena. El estudio debe ser explícitamente claro en que \textit{no
constituye diagnóstico}, y debe evitar el uso exclusivo de terminología
anglohablante (TEA, 2e, \textit{masking}) que puede generar resistencia o
incomprensión. El uso de narrativas locales y culturalmente situadas
(``personas que piensan diferente y tienen talentos especiales'') puede
facilitar la participación sin trivializar el fenómeno.

\subsubsection{Estigma en contexto costeño}

La cultura regional de la Costa Caribe colombiana tiene dinámicas específicas
en relación al estigma psiquiátrico que deben tenerse en cuenta en el diseño
del estudio. La identidad de género masculina tradicional, predominante en la
planta docente histórica de Unisucre, puede generar inhibición frente a reconocer
dificultades sociales o emocionales. La estrategia de enmarcar el estudio como
una investigación sobre ``perfiles de excelencia y neurodiversidad'' en lugar de
sobre ``trastornos'' puede reducir esta barrera.

\subsubsection{Potencial de integración con formación de posgrado}

Unisucre cuenta con maestrías en Ciencias de la Educación y en Biología.
Los docentes estudiantes de posgrado (una población con alta representación
de perfiles 2e según la evidencia internacional) representan una subpoblación
especialmente pertinente. Adicionalmente, la integración de esta investigación
como proyecto de trabajo de grado de maestría permitiría multiplicar la capacidad
investigadora disponible sin requerir financiación externa significativa.

\subsubsection{Articulación con MINCIENCIAS y agenda de ciencia abierta}

La investigación propuesta puede articularse con las líneas de financiación de
MINCIENCIAS en educación de calidad (ODS 4) y reducción de desigualdades (ODS 10).
La invisibilidad de académicos con TEA y AACC en el sistema científico colombiano
es, en última instancia, un problema de \textit{equidad en la producción del
conocimiento}: si las barreras institucionales para los perfiles neurodivergentes
determinan qué tipo de investigadores pueden sostenerse en el sistema, también
determinan qué tipo de conocimiento se produce y qué perspectivas quedan
sistemáticamente excluidas.

\subsection{Posibles hallazgos y contribuciones anticipadas}

Aunque no es posible predeterminar los resultados de una investigación rigurosa,
la extrapolación cuidadosa de la literatura existente permite anticipar que un
estudio en Unisucre podría documentar:

\begin{enumerate}[leftmargin=*, itemsep=3pt]
    \item \textbf{Una prevalencia de rasgos 2e superior a la estimada}, dado el
    subdiagnóstico estructural del contexto colombiano y la mayor representación
    de rasgos autísticos en entornos académicos.
    \item \textbf{Estrategias de enmascaramiento altamente desarrolladas}, con
    consecuencias en salud mental (burnout, ansiedad, sensación crónica de
    impostura) raramente conectadas por los afectados con su perfil neurodivergente.
    \item \textbf{Ausencia de diagnóstico formal en la mayoría de los casos},
    con autoidentificación reciente facilitada por el acceso a información en
    redes sociales y plataformas digitales —un fenómeno documentado en otros
    contextos latinoamericanos \parencite{pessoa2026}.
    \item \textbf{Fortalezas investigadoras específicas} asociadas al perfil 2e
    (pensamiento divergente, atención sostenida, profundidad de análisis,
    pensamiento sistémico no convencional) que el sistema académico actual
    solo aprovecha parcialmente o activamente penaliza.
    \item \textbf{Necesidades de apoyo institucional simples y de bajo coste}
    que no requieren grandes inversiones pero sí voluntad política:
    comunicación escrita de instrucciones, agendas predecibles, posibilidad de
    trabajo en entornos tranquilos, revisión de los sistemas de evaluación docente.
\end{enumerate}

% ─────────────────────────────────────────────────────────────
%  ENTRADAS ADICIONALES para referencias.bib
%  Añadir al final del archivo referencias.bib existente
% ─────────────────────────────────────────────────────────────

% @article{foleynicpon2011_paradox,
%   author  = {Foley-Nicpon, Megan and Allmon, Alissa and Sieck, Brigit and
%              Stinson, Renee D.},
%   title   = {Empirical Investigation of Twice-Exceptionality: Where Have We
%              Been and Where Are We Going?},
%   journal = {Gifted Child Quarterly},
%   year    = {2011},
%   volume  = {55},
%   number  = {1},
%   pages   = {3--17},
%   doi     = {10.1177/0016986210382575}
% }
% 
% @article{cain2019,
%   author  = {Cain, Meghan K. and Kaboski, Juhi R. and Gilger, Jeffrey W.},
%   title   = {Profiles and Academic Trajectories of Cognitively Gifted
%              Children with Autism Spectrum Disorder},
%   journal = {Autism},
%   year    = {2019},
%   volume  = {23},
%   number  = {7},
%   pages   = {1663--1674},
%   doi     = {10.1177/1362361318804019}
% }
% 
% @article{assouline2021,
%   author  = {Assouline, Susan G. and Foley-Nicpon, Megan},
%   title   = {Processing Speed as a Mediator of Academic Achievement in
%              Gifted Youth with Autism Spectrum Disorder},
%   journal = {Journal of Autism and Developmental Disorders},
%   year    = {2021},
%   doi     = {10.1007/s10803-021-05290-4}
% }
% 
% @article{reis2021,
%   author  = {Reis, Sally M. and Gelbar, Nicholas W. and Madaus, Joseph W.},
%   title   = {Understanding the Academic Success of Academically Talented
%              College Students with Autism Spectrum Disorders},
%   journal = {Journal of Autism and Developmental Disorders},
%   year    = {2021},
%   doi     = {10.1007/s10803-021-05290-4}
% }
% 
% @article{madaus2022factors,
%   author  = {Madaus, Joseph and Reis, Sally and Gelbar, Nicholas and
%              Delgado, Julie and Cascio, Alexandra},
%   title   = {Perceptions of Factors That Facilitate and Impede Learning
%              Among Twice-Exceptional College Students with Autism Spectrum
%              Disorder},
%   journal = {Neurobiology of Learning and Memory},
%   year    = {2022},
%   doi     = {10.1016/j.nlm.2022.107627}
% }
% 
% @article{madaus2022transition,
%   author  = {Madaus, Joseph and Cascio, Alexandra and Delgado, Julie and
%              Gelbar, Nicholas and Reis, Sally and Tarconish, Emily},
%   title   = {Improving the Transition to College for Twice-Exceptional
%              Students with {ASD}: Perspectives From College Service Providers},
%   journal = {Career Development and Transition for Exceptional Individuals},
%   year    = {2022},
%   doi     = {10.1177/21651434221091230}
% }
% 
% @article{austermann2023,
%   author  = {Austermann, Quinn and Gelbar, Nicholas W. and Reis, Sally M.
%              and Madaus, Joseph W.},
%   title   = {The Transition to College: Lived Experiences of Academically
%              Talented Students with Autism},
%   journal = {Frontiers in Psychiatry},
%   year    = {2023},
%   doi     = {10.3389/fpsyt.2023.1125904}
% }
% 
% @article{macleod2017,
%   author  = {MacLeod, Andrea and Allan, Julie and Lewis, Ann
%              and Robertson, Chris},
%   title   = {`Here I Come Again': The Cost of Success for Higher Education
%              Students Diagnosed with Autism},
%   journal = {International Journal of Inclusive Education},
%   year    = {2017},
%   volume  = {21},
%   number  = {11},
%   pages   = {1165--1177},
%   doi     = {10.1080/13603116.2017.1396502}
% }
% 
% @article{gillespie-lynch2017,
%   author  = {Gillespie-Lynch, Kristen and Bublitz, Dennis and Donachie, Annemarie
%              and Wong, Vincent and Brooks, Patricia J. and D'Onofrio, Joanne},
%   title   = {``For a Long Time Our Voices Have Been Hushed'': Using Student
%              Perspectives to Develop Supports for Neurodiverse College Students},
%   journal = {Frontiers in Psychology},
%   year    = {2017},
%   volume  = {8},
%   pages   = {544},
%   doi     = {10.3389/fpsyg.2017.00544}
% }
% 
% @article{perkowski2017,
%   author  = {Perkowski, Maciej and Oksztulski, Maciej and Kaczyńska, Izabela},
%   title   = {Autism Between the {PhD} Student and the Promotor: A Case Study},
%   journal = {Studies in Logic, Grammar and Rhetoric},
%   year    = {2017},
%   doi     = {10.1515/slgr-2017-0048}
% }
% 
% @article{wallpolin2017,
%   author  = {Wall-Polin, R. K.},
%   title   = {Relative to {Einstein}: Quality of Life in Twice Exceptional Adults},
%   year    = {2017},
%   doi     = {10.33015/dominican.edu/2017.psy.rp.03},
%   note    = {Dominican University of California, Research Posters}
% }
% 
% @article{alhroub2021,
%   author  = {Al-Hroub, Anies},
%   title   = {Utility of Psychometric and Dynamic Assessments for Identifying
%              Cognitive Characteristics of Twice-Exceptional Students},
%   journal = {Frontiers in Psychology},
%   year    = {2021},
%   volume  = {12},
%   pages   = {747872},
%   doi     = {10.3389/fpsyg.2021.747872}
% }
% 
% @article{burgerveltmeijer2023,
%   author  = {Burger-Veltmeijer, Agnes and Minnaert, Alexander},
%   title   = {Needs-Based Assessment of Twice-Exceptional Gifted Students:
%              The {S\&W-Heuristic}},
%   journal = {Advances in Social Sciences Research Journal},
%   year    = {2023},
%   doi     = {10.14738/assrj.101.13830}
% }
% 
% @article{linares2026,
%   author  = {Liñares-de-Marcos, Jon and Palomero-Sierra, Blanca and
%              Sánchez-Gómez, Victoria and Fernández-Álvarez, Clara J.
%              and others},
%   title   = {An Integrative Approach Between Neurodiversity Perspectives and
%              Quality of Life Models for Autistic People Across the Spectrum of
%              Support Needs},
%   journal = {Frontiers in Psychiatry},
%   year    = {2026},
%   doi     = {10.3389/fpsyt.2025.1756323}
% }
% 
% @article{pessoa2026,
%   author  = {Pessoa de Souza, Carla and Radek, Ereni},
%   title   = {High Abilities and Giftedness in Adolescents and Adults with Autism},
%   journal = {Revista Gênero e Interdisciplinaridade},
%   year    = {2026},
%   doi     = {10.51249/gei.v6i06.2781}
% }
% 
% @article{poirier2025,
%   author  = {Poirier, Joanie and Brault-Labbé, Anne and Brassard, Audrey},
%   title   = {Living With the Gift of Giftedness: An Exploratory Study on
%              the Well-Being of Intellectually Gifted Adults},
%   journal = {Gifted Child Quarterly},
%   year    = {2025},
%   doi     = {10.1177/00169862251347293}
% }
% 
% @book{smith2009,
%   author    = {Smith, Jonathan A. and Flowers, Paul and Larkin, Michael},
%   title     = {Interpretative Phenomenological Analysis: Theory, Method and
%                Research},
%   year      = {2009},
%   publisher = {SAGE},
%   address   = {London}
% }
% 
% @book{yin2018,
%   author    = {Yin, Robert K.},
%   title     = {Case Study Research and Applications: Design and Methods},
%   edition   = {6th},
%   year      = {2018},
%   publisher = {SAGE},
%   address   = {Thousand Oaks, CA}
% }
% 
% @book{kemmis2014,
%   author    = {Kemmis, Stephen and McTaggart, Robin and Nixon, Rhonda},
%   title     = {The Action Research Planner: Doing Critical Participatory
%                Action Research},
%   year      = {2014},
%   publisher = {Springer},
%   address   = {Singapore}
% }
% 
% @article{gonzalez2010,
%   author  = {González Fontao, María del Pilar and Martínez Suárez, Montserrat},
%   title   = {Identificación del alumnado con altas capacidades intelectuales},
%   journal = {Revista Educación Inclusiva},
%   year    = {2010},
%   volume  = {3},
%   number  = {1},
%   pages   = {131--143}
% }
% 
% @article{szabó2018,
%   author  = {Szabó, János and Révész, György},
%   title   = {Eternal Questions of Gifted Education from the Aspect of
%              University Teachers},
%   journal = {Journal for the Education of Gifted Young Scientists},
%   year    = {2018},
%   doi     = {10.17478/jegys.2018.72}
% }
% 
% @article{reuwsaat2026,
%   author  = {Reuwsaat, Karin and Poltronieri, Bruno and Panizzutti, Rogério
%              and Velasques, Bruna},
%   title   = {Giftedness: A Critical Analysis of Theories and Identification
%              Methods in Light of Contemporary Neuroscience},
%   journal = {International Journal of Developmental Neuroscience},
%   year    = {2026},
%   doi     = {10.1002/jdn.70126}
% }
